\documentclass[10pt,a4paper]{amsart}
\usepackage{amsmath,amssymb,amsthm,graphicx}
% For commutative diagrams we use Xy-pic
\usepackage[all]{xy}

% We use 2cell for 2-commutative diagrams.
\xyoption{2cell}
% Theorem environments.
%
\theoremstyle{plain}
\newtheorem{theorem}[subsection]{Theorem}
\newtheorem{proposition}[subsection]{Proposition}
\newtheorem{lemma}[subsection]{Lemma}

\theoremstyle{definition}
\newtheorem{definition}[subsection]{Definition}
\newtheorem{example}[subsection]{Example}
\newtheorem{exercise}[subsection]{Exercise}
\newtheorem{situation}[subsection]{Situation}

\theoremstyle{remark}
\newtheorem{remark}[subsection]{Remark}
\newtheorem{remarks}[subsection]{Remarks}

\numberwithin{equation}{subsection}

% Macros
%
\def\lim{\mathop{\mathrm{lim}}\nolimits}
\def\colim{\mathop{\mathrm{colim}}\nolimits}
\def\Spec{\mathop{\mathrm{Spec}}}
\def\Hom{\mathop{\mathrm{Hom}}\nolimits}
\def\Ext{\mathop{\mathrm{Ext}}\nolimits}
\def\SheafHom{\mathop{\mathcal{H}\!\mathit{om}}\nolimits}
\def\SheafExt{\mathop{\mathcal{E}\!\mathit{xt}}\nolimits}
\def\Sch{\mathit{Sch}}
\def\Mor{\mathop{\mathrm{Mor}}\nolimits}
\def\Ob{\mathop{\mathrm{Ob}}\nolimits}
\def\Sh{\mathop{\mathit{Sh}}\nolimits}
\def\NL{\mathop{N\!L}\nolimits}
\def\CH{\mathop{\mathrm{CH}}\nolimits}
\def\proetale{{pro\text{-}\acute{e}tale}}
\def\etale{{\acute{e}tale}}
\def\QCoh{\mathit{QCoh}}
\def\Ker{\mathop{\mathrm{Ker}}}
\def\Im{\mathop{\mathrm{Im}}}
\def\Coker{\mathop{\mathrm{Coker}}}
\def\Coim{\mathop{\mathrm{Coim}}}

% Boxtimes
%
\DeclareMathSymbol{\boxtimes}{\mathbin}{AMSa}{"02}

%
% Macros for moduli stacks/spaces
%
\def\QCohstack{\mathcal{QC}\!\mathit{oh}}
\def\Cohstack{\mathcal{C}\!\mathit{oh}}
\def\Spacesstack{\mathcal{S}\!\mathit{paces}}
\def\Quotfunctor{\mathrm{Quot}}
\def\Hilbfunctor{\mathrm{Hilb}}
\def\Curvesstack{\mathcal{C}\!\mathit{urves}}
\def\Polarizedstack{\mathcal{P}\!\mathit{olarized}}
\def\Complexesstack{\mathcal{C}\!\mathit{omplexes}}
% \Pic is the operator that assigns to X its picard group, usage \Pic(X)
% \Picardstack_{X/B} denotes the Picard stack of X over B
% \Picardfunctor_{X/B} denotes the Picard functor of X over B
\def\Pic{\mathop{\mathrm{Pic}}\nolimits}
\def\Picardstack{\mathcal{P}\!\mathit{ic}}
\def\Picardfunctor{\mathrm{Pic}}
\def\Deformationcategory{\mathcal{D}\!\mathit{ef}}
\newenvironment{reference}{}{}
\newenvironment{slogan}{}{}
\newenvironment{history}{}{}
\usepackage{cleveref}
\crefname{originallemma}{original lemma}{original lemmas}
\Crefname{originallemma}{Original lemma}{Original lemmas}
\makeatletter
\newlabel{lemma-regular-model-connected}{{0C62}{source}{}{sourcepointer-1}{}}
\newlabel{lemma-regular-model-connected@cref}{{[originallemma][0][]0C62}{[1][0][]source}{}{}{}}
\makeatother
\crefname{originallemma}{original lemma}{original lemmas}
\Crefname{originallemma}{Original lemma}{Original lemmas}
\makeatletter
\newlabel{lemma-minimal-model-unique}{{0C6B}{source}{}{sourcepointer-2}{}}
\newlabel{lemma-minimal-model-unique@cref}{{[originallemma][0][]0C6B}{[1][0][]source}{}{}{}}
\makeatother
\crefname{originallemma}{original lemma}{original lemmas}
\Crefname{originallemma}{Original lemma}{Original lemmas}
\makeatletter
\newlabel{lemma-nonuniqueness}{{0CDA}{source}{}{sourcepointer-3}{}}
\newlabel{lemma-nonuniqueness@cref}{{[originallemma][0][]0CDA}{[1][0][]source}{}{}{}}
\makeatother
\title{A smooth proper model, a smooth minimal model and smoothness of every minimal model are equivalent}
\author{The Stacks Project Authors}
\date{stacks-project:models.tex:local-snapshot-20261012; GNU FDL 1.2 or later}
\begin{document}
\maketitle
\noindent\textit{Editorial note.} The counted unit is the Stacks Project Lemma lemma-good at source lines 6415-6425 together with its proof at source lines 6427-6439, a straight adjacent pairing in the chapter file. The chapter is the upstream file models.tex, which is the file name recorded in the source record; the pool directory carries the longer name stacks-models and is not used as the chapter name anywhere in this card.\par\medskip
\noindent\textit{Editorial note.} This chapter's candidate index records the statement body and the proof body without their closing markers, so the retained spans used here are the full environments, exactly as in the earlier cards of this chapter.\par\medskip
\noindent\textit{Editorial note.} The retained proof names three results of the same chapter, one for the connectedness of the special fibre of a regular model, one for the uniqueness of the minimal model in positive genus and one for the structure of the special fibre when the minimal model is not unique. All three are declared as non-load-bearing navigation pointers with their stable Stacks tags looked up in the shared tag index; the counted argument is the assembly of the three implications.\par\medskip
\noindent\textit{Editorial note.} The statement mentions no situation or labelled object, so no context span is retained and no bibliography entry is needed.\par\medskip
\noindent\textit{Editorial note.} The retained preamble is the chapter preamble's Xy-pic setup at lines 14-18, the theorem declarations and mathematical macros at lines 154-217, and the three no-op wrapper environments.\par\medskip
\noindent\textit{Allowed standard prerequisites (editorial).}\begin{itemize}
\item Models of a curve over a discrete valuation ring and smoothness of the special fibre. A model of a smooth projective curve over the fraction field is a flat proper scheme over the discrete valuation ring with that curve as its generic fibre, and it is called minimal when it admits no contraction, equivalently when its numerical type is minimal; the special fibre of a regular model with connected components is connected, and a smooth proper model has smooth connected special fibre. The author's own definitions and the neighbouring lemmas of the same chapter are the allowed inputs.. The counted lemma collects the three equivalent formulations of good reduction that the chapter uses when it states its main results on semistable reduction.
\end{itemize}\medskip
\ifdefined\hypertarget\hypertarget{sourcepointer-1}{}\fi
\noindent\textit{Source reference.} Original lemma 0C62: focused-seeds/AGP/sources/stacks-snapshot-20261011/tags.txt, source lines 15282--15282. Navigation pointer invoked by the retained proof for the connectedness of the special fibre; the cited statement is an allowed input and is not reproduced inside the span. The number is the stable Stacks tag of the target, looked up in the shared tag index.\par\medskip
\ifdefined\hypertarget\hypertarget{sourcepointer-2}{}\fi
\noindent\textit{Source reference.} Original lemma 0C6B: focused-seeds/AGP/sources/stacks-snapshot-20261011/tags.txt, source lines 15291--15291. Navigation pointer invoked by the retained proof for the uniqueness of the minimal model in positive genus; that uniqueness is an allowed input and is not reproduced inside the span. The number is the stable Stacks tag of the target, looked up in the shared tag index.\par\medskip
\ifdefined\hypertarget\hypertarget{sourcepointer-3}{}\fi
\noindent\textit{Source reference.} Original lemma 0CDA: focused-seeds/AGP/sources/stacks-snapshot-20261011/tags.txt, source lines 15535--15535. Navigation pointer invoked by the retained proof in the case where the minimal model is not unique; the cited analysis of the special fibre is an allowed input and is not reproduced inside the span. The number is the stable Stacks tag of the target, looked up in the shared tag index.\par\medskip
\begin{lemma}
\label{lemma-good}
Let $R$ be a discrete valuation ring with fraction field $K$.
Let $C$ be a smooth projective curve over $K$ with $H^0(C, \mathcal{O}_C) = K$.
The following are equivalent
\begin{enumerate}
\item there exists a proper smooth model for $C$,
\item there exists a minimal model for $C$ which is smooth over $R$,
\item any minimal model is smooth over $R$.
\end{enumerate}
\end{lemma}

\begin{proof}
If $X$ is a smooth proper model, then the special fibre is
connected (Lemma \ref{lemma-regular-model-connected})
and smooth, hence irreducible. This immediately implies that
it is minimal. Thus (1) implies (2). 
To finish the proof we have to show that (2) implies (3).
This is clear if the genus of $C$ is $> 0$, since then
the minimal model is unique (Lemma \ref{lemma-minimal-model-unique}).
On the other hand, if the minimal model is not unique, then
the morphism $X \to \Spec(R)$ is smooth for any minimal model
as its special fibre will be isomorphic to $\mathbf{P}^1_k$
by Lemma \ref{lemma-nonuniqueness}.
\end{proof}

\end{document}
